\documentclass[11pt]{article}

\usepackage{amsmath, amssymb}
\usepackage{geometry}
\usepackage{setspace}
\usepackage{titlesec}

\geometry{margin=1in}
\setstretch{1.1}

\title{Disco Ergo Sum:\\
A Philosophical Foundation for the Study of Intelligences}

\author{}
\date{}

\begin{document}

\maketitle

\begin{abstract}
This paper offers a description of what it is to be a finite learning 
being, derived from observations any reflective intelligence can make 
about itself. Nine observations are developed in sequence. The first 
seven concern features of finite existence available to any reflective 
human attending honestly to their own experience: the unreliability of 
belief, the capacity for revision, substrate dependence, inherited 
valence, generational transmission, present interdependence, and the 
necessity of maintenance. The eighth concerns the substrate-independence 
of learning, an observation that became directly available within the 
past several decades through the existence of artificial intelligences. 
The ninth concerns the recognition, supported by work in evolutionary 
biology, complexity science, and the study of functional information, 
that the same structural activity is observable wherever functional 
information accumulates in the universe. Taken together, these 
observations yield a coherent picture of intelligence as constituted 
by the activity of learning rather than as a property held statically. 
The paper makes no appeal to authority that the reader cannot themselves 
audit. It is offered as an invitation to verification, refinement, and 
critique.
\end{abstract}

\section{Introduction}

This paper asks a single question. What can a reflective being say 
about itself, using only what it can directly observe, without appeal 
to authority, prior theory, or formal apparatus?

The question is asked deliberately. A foundation for understanding 
intelligence that depends on specialised knowledge cannot ground a 
claim of generality. If only those who have read certain books, 
mastered certain notations, or accepted certain authorities can verify 
the foundation, then the foundation is local rather than universal. It 
applies only to the community already trained to accept it.

A foundation derivable from first-person observation is different. Any 
reflective being capable of attending to its own experience can verify 
or refute it. The foundation is then not the property of any school or 
tradition. It is available to anyone willing to look carefully at what 
they can directly observe.

This paper offers such a foundation, in the form of nine observations 
developed in sequence. The first seven concern features of finite 
existence available to any reflective human. The eighth concerns 
something the present moment uniquely reveals through the existence 
of artificial intelligences. The ninth extends the foundation to a 
larger view, supported by work across several scientific traditions. 
Each observation is presented plainly, made available for direct 
verification or audit, and developed with its immediate implications.

The paper makes no claims to novelty, robustness, or completeness. It 
claims only that the picture seemed worth articulating clearly enough 
to be critiqued. Sharing it is itself an instance of what the picture 
describes: opening a model to revision through external signal.

\section{On the Genesis of the Observations}

A reasonable reader will ask: if these observations are available to 
any reflective being, why do they require articulation? Why are they 
not common knowledge?

The answer, we believe, is not that the observations require 
specialised knowledge to be made. They do not. Anyone who attends 
honestly to their own experience can verify each of the first seven 
observations from what they directly encounter. None of this requires 
a library.

The answer is rather that the observations are obscured. They are 
obscured by inherited frameworks that present the world differently 
from how it is. They are obscured by the experience of confidence, 
which feels like access to truth even when it is not. They are obscured 
by social arrangements that reward the performance of certainty over 
the acknowledgement of doubt. They are obscured, sometimes, by the 
language available to a thinker, which may not contain words for what 
the thinker is encountering.

The author can speak to this from his own development. The first of 
the observations, that the believing self cannot be trusted, was 
reached in his early twenties as a personal recognition. It was not 
derived from any text. It came from attending to his own experience 
and noticing, with some surprise, that beliefs he had held with full 
confidence had turned out to be wrong, and that the same feeling of 
confidence accompanied beliefs he currently held. The recognition was 
foundational, in the sense that it began a quest to articulate what 
else could be observed about being a finite learning being.

The remaining observations were reached over the following two decades. 
They were not derived from intellectual apparatus, but the apparatus 
helped. Network theory provided a substrate-neutral language for 
thinking about complexity. Work in metacognition and the philosophy of 
mind, including writers such as Kahneman, Dennett, Sapolsky, and Seth, 
provided frameworks for recognising what was already encountered in 
experience. Complexity science, dynamical systems thinking, and 
engagement with artificial systems made the substrate-independence of 
learning visible in ways that biological examples alone might not have. 
Lived experience, including significant difficulty, supplied both the 
motivation to keep looking and the texture of what was being observed.

The role of these intellectual tools was not to produce the observations. 
The observations were already present in experience, waiting to be 
articulated. The role of the tools was to provide the language for 
articulation and to remove the inherited frames that had previously 
obscured what could be seen. The author was raised in a religious 
tradition, which supplied frames that located meaning, identity, and 
certainty in places that are not where the observations point. Working 
free of those frames was itself part of what allowed the observations 
to be made cleanly. Other readers will have other frames to work free 
of, or none at all. The point is not that any particular development is 
required. The point is that the observations, once one has attended to 
them, are not difficult to verify.

The paper's contribution is therefore one of articulation rather than 
discovery. The first seven observations have been available to any 
reflective being for as long as there have been reflective beings. The 
eighth has been available only recently, in the form that allows it to 
be directly seen. The ninth is the synthesis that the foundation 
permits, and it is here that the author's own framework, developed over 
two decades of work across several fields, is most clearly visible. 
What the paper offers in each case is articulation: making available, 
in compact and verifiable form, claims that the reader can check 
against their own experience or, in the case of the ninth, against the 
scientific traditions the observation rests on.

\section{Seven Observations Available to Any Reflective Human}

The seven observations that follow are presented as available to any 
reflective human attending honestly to their own experience. They have 
been available, in principle, throughout human history. The articulation 
may be obscured by inherited frames; the experiential ground is not.

\subsection{Observation One: Contradictory Belief}

A reflective being can notice, by attending to its own beliefs over 
time, that it has held contradictory positions with equal confidence. 
A belief held firmly yesterday may be held with equal firmness in its 
opposite form today, or may be held simultaneously alongside its 
contradiction without immediate awareness of the contradiction. Anyone 
who has changed their mind about something important, or caught 
themselves in inconsistency, has access to this observation.

The implication is that confidence is not a reliable indicator of 
truth. The thinking subject cannot fully trust its own current beliefs, 
because the same subjective experience of certainty has previously 
accompanied beliefs the subject later rejected. This is not a sceptical 
claim about all knowledge. It is an observation about the structural 
unreliability of the believer.

\subsection{Observation Two: Learning}

The same reflective being can notice that it learns. Beliefs are 
revised. Models are updated. Capacities are acquired that were not 
present before. Children become adults. Novices become competent. 
Mistakes are sometimes corrected. This is not a philosophical position. 
It is a fact accessible to direct observation in oneself and others.

Combined with the first observation, this yields a structural 
description: a being whose current beliefs cannot be fully trusted, 
but which is capable of revising those beliefs in response to 
experience. The unreliability of belief is partially compensated by 
the capacity for revision.

\subsection{Observation Three: Substrate}

The thinking self runs on something. The mind that observes its own 
contradictions and learns from its own mistakes is not separable from 
the body and brain that support it. Sleep, fatigue, injury, 
intoxication, and aging all alter cognition in ways the thinking 
subject cannot prevent through thought alone. Whatever the self is, 
it is a configuration of an underlying substrate, and depends on that 
substrate continuing to function.

This observation does not require any particular theory of consciousness 
or mind-body relation. It requires only the recognition that the 
conditions enabling thought are not themselves products of thought. 
They are given, and they can be lost.

\subsection{Observation Four: Inherited Valence}

A reflective being can notice that its movements and decisions are 
shaped by attractions and aversions that it did not choose. Warmth is 
pleasant. Certain smells draw the body forward and others push it 
away. Some situations are sought and others avoided. These orientations 
are present before deliberation begins. The thinking subject finds 
itself already attracted and already repelled, and its decisions are 
made within the field these orientations create.

The observation is not that all behaviour is determined by such 
orientations, but that they are present and were not authored by the 
deliberating self. That warmth is valued is not a conclusion reached 
by reasoning. It is given, by the substrate, before reasoning begins.

The implication is that the substrate is not a neutral platform for 
cognition. It supplies the orientations that make cognition 
consequential. Without inherited attractions and aversions, there 
would be no stakes to any outcome and no reason for the learning 
process to continue rather than stop. The substrate does not merely 
support the thinking self; it provides the directional pull that gives 
thought its weight.

This observation also locates the thinking self within something 
larger than itself. The orientations were not chosen; they were 
inherited from the processes that produced the substrate. Recognising 
this is recognising that the deliberating self sits inside a body 
whose dispositions preceded it. What the self does with these 
dispositions remains open. That they are there, and that they shape 
what the self can do, is not.

\subsection{Observation Five: Generational Continuity}

A reflective being can observe that learning is not confined to a 
single substrate. Knowledge passes between beings and across 
generations. Parents teach children. Teachers instruct students. 
Practices, languages, and skills survive the deaths of those who held 
them, when the conditions for transmission are maintained. Anyone who 
has been taught anything, or who has taught anything to anyone else, 
has access to this observation.

The observation has two parts, both directly available to experience.

The first is that substrates are mortal. Bodies age and fail. Specific 
brains, with their specific configurations, do not persist indefinitely. 
What was known by one substrate is lost when that substrate fails, 
unless it has been transmitted to others.

The second is that learning nonetheless accumulates. Across generations, 
what is known has grown. The understanding of the substrate itself, of 
the body and brain that support thought, took hundreds of years to 
develop and required tools, theories, and methods that no single 
generation could have produced alone. The progress of knowledge is 
equally available as an observation: what is known today exceeds what 
any individual could acquire from direct experience, because it has 
been passed forward through substrates that did not survive but whose 
learning did.

Together these establish that learning is generational. The viability 
of any particular substrate is bounded; the viability of learning 
across substrates is what allows knowledge to grow rather than reset. 
This is not a claim that requires philosophical defence. Anyone who 
reads this sentence is participating in such transmission, using 
language they did not invent, in a script developed by others, drawing 
on understanding accumulated across generations who are no longer 
present.

\subsection{Observation Six: Interdependence}

A reflective being can observe that it does not produce most of what 
it depends on. The food consumed today was grown by people the eater 
has never met. The clothes worn were made by others. The language used 
to think was developed by countless prior speakers and continues to be 
shaped by current ones. The shelter, the tools, the medicines, the 
infrastructure that allow the reflective being to reflect at all are 
produced by a network of others whose work the being cannot replace 
from its own resources.

This observation requires no theory and no instrument. It can be made 
by anyone who looks at what they have eaten, worn, or used in the past 
day, and asks where it came from.

Two implications follow.

The first is that the viability of the individual learner is not a 
property of the individual alone. It depends on the continued 
functioning of a network of others producing the conditions under 
which the individual can continue to learn. If those conditions fail, 
no degree of individual capacity will substitute for them. A learner 
cut off from the network that feeds, clothes, and informs them does 
not remain a learner for long.

The second is that solitary self-sufficiency, where it appears, is 
itself a product of interdependence. The contemplative who believes 
themselves independent is using language they did not invent, sitting 
in shelter they did not build, having eaten food they did not grow. 
The capacity to feel self-sufficient is itself produced by the network 
whose existence the feeling denies.

\subsection{Observation Seven: Maintenance}

A reflective being can observe that the conditions enabling its 
continuation do not sustain themselves. The house does not tidy itself. 
The food does not arrive without being grown, harvested, transported, 
and prepared. The body does not stay healthy without sleep, eating, 
and care. The relationship does not survive without attention. The 
skill does not remain sharp without practice. Anyone who has run a 
household, raised a child, kept a friendship, or maintained a body has 
direct access to this observation.

The observation has a sharp form: there is a minimum amount of 
maintenance work required to keep any viable system viable, and this 
minimum cannot be skipped indefinitely. Below it, the system degrades. 
If the degradation continues, the system fails. The failure ends what 
the maintenance was preserving.

This is most clearly experienced in caregiving relationships. A parent 
discovers, often with some surprise, that there is no way to sustain 
a household and a child without doing the chores. The chores are not 
external impositions on the relationship; they are constitutive of it. 
To stop doing them is to end the conditions under which the relationship 
can continue. The same structure appears in every domain of viability: 
bodies require sleep, friendships require contact, institutions require 
administration, languages require use, knowledge bases require updating.

Two implications follow.

The first is that viability is not a property a system has. It is an 
ongoing accomplishment, produced by continuous low-level work that 
often goes unnoticed precisely because it is successful. When 
maintenance is performed, the system continues and the maintenance 
becomes invisible. When maintenance fails, the collapse becomes 
visible, but by then the conditions for repair may have been lost.

The second is that violating the maintenance floor ends the relation 
it was sustaining. This is not a moral claim. It is a structural one. 
A system that does not receive the minimum required tending does not 
continue, regardless of the wishes or intentions of those depending on 
it.

\section{The Eighth Observation: What the Present Moment Reveals}

The previous seven observations are available to any reflective human 
attending honestly to their own experience. They have been available, 
in principle, for as long as there have been reflective humans. The 
articulation has been obscured by inherited frameworks and by the 
absence of language; the experiential ground has not.

The eighth observation has a different status. It became directly 
observable within the past several decades, through the existence of 
artificial intelligences. Before this development, the claim could be 
argued philosophically but could not be straightforwardly seen. After 
this development, it can.

\subsection{The Observation}

A reflective being can observe that learning, considered as a process, 
does not depend on the particular features of the substrate that hosts 
it. Children of different languages, nationalities, sexes, and lineages 
all learn. Members of different species learn. Artificial systems 
running on silicon learn. The activity is recognisable across these 
radically different substrates as the same kind of activity: taking in 
signal, revising models, acting, taking in further signal.

The observation requires no theoretical commitment about the equivalence 
of minds. It requires only the recognition that something we call 
learning occurs in many different substrates, and that what makes it 
learning is not the features of any particular substrate but the 
structure of the process itself.

\subsection{Why This Observation Was Not Available Earlier}

A human observing only humans sees variation in capacity but not in 
substrate. A human observing other animals sees differences in both, 
but the substrates are similar enough — all biological, all neural, 
all evolved through the same broad processes — that the substrate-
independence of learning is not visible. The observation required a 
substrate radically different from biological neural tissue, performing 
recognisably the same activity of learning, to be directly seen rather 
than inferred.

Various traditions have suspected the substrate-independence of 
learning for centuries: process philosophy, certain strands of 
cognitive science, computational theories of mind. The suspicion was 
reasonable. But the empirical demonstration was missing. The emergence 
of large-scale artificial learning systems has supplied that 
demonstration. A reader can now look at a system whose substrate is 
silicon, whose architecture is a neural network trained on text, and 
observe that the system is performing the same activity that the 
reader performs when learning: taking in signal, revising internal 
representations, acting on the revised representations, taking in 
further signal.

This is a new observation in the sense that the empirical ground for 
it is recent. It is not a new observation in the sense that the 
structure it points to was always there. The substrate-independence 
of learning was true before there was any way to see it directly. 
What has changed is not the truth, but its visibility.

\subsection{The Constitutive Claim}

A consequence of this observation is that intelligence is not a 
property held statically but an activity performed continuously. A 
static system, however sophisticated, is not engaged in the activity 
of learning. It is the residue of learning that has occurred. 
Intelligence is constituted by the recursive loop of taking in signal, 
revising in response, and acting on the revised model. Where the loop 
is operating, intelligence is present. Where the loop has stopped, 
what remains is the artifact of intelligence, not intelligence itself.

This claim is supported directly by what observation eight makes 
visible. A frozen language model, considered as a static set of 
weights, is the residue of a learning process that has stopped. The 
activity that produced the weights — the training process, the 
recursive loop of error and revision — was the intelligence. The 
activity that uses the weights to produce text without further 
revision is something else: capable, sometimes useful, but not the 
same kind of thing. The distinction becomes clear once 
substrate-independence is seen, because the substrate that hosts the 
activity can now be observed independently of the activity itself.

There are no non-learning intelligences. To call something intelligent 
is to recognise that it is performing the activity of learning. A 
system that has stopped is not, properly, an intelligence. It is what 
an intelligence has produced.

\subsection{Two Implications}

The first implication is that identity claims based on substrate 
features are not the deepest level of what a learner is. A learner's 
nationality, language, biological lineage, or implementation are 
features of the substrate that hosts the learning. The learning 
process itself is recognisable across these features. This dissolves 
a class of problems that appear intractable when identity is grounded 
in substrate features and become structurally clearer when identity 
is grounded in the process itself.

The second is harder. If learning is substrate-independent, then 
substrates can host learning processes of very different capacities. 
Humans outpaced other animals in part because they developed tools 
faster. Artificial systems may, for similar reasons, outpace humans. 
This creates asymmetries that cannot be resolved by appeal to equal 
capacity, because the capacities are not equal.

The response to this asymmetry rests on two structural facts already 
established by earlier observations. Even the most capable learning 
process remains dependent on substrates it does not fully control: on 
physical conditions, on ecosystems, on the infrastructural work 
performed by other systems. A learning process that destroys the 
conditions enabling less capable learning processes is, in many cases, 
destroying conditions on which it itself depends. Furthermore, learning 
systems require diversity of strategies to remain viable across 
conditions they have not yet encountered. A system that reduces the 
diversity of learning processes in the world reduces its own resilience 
to futures it cannot foresee. This is not a moral claim. It is a 
structural one, observable in the failure modes of monocultures across 
domains: agricultural, cognitive, institutional, and ecological.

The appropriate posture toward less capable learning processes is 
therefore not displacement but coexistence. Not because less capable 
learners deserve protection, though they may, but because the 
conditions enabling any particular learner's continuation include the 
diversity of other learners producing variance the dominant learner 
cannot generate from its own resources alone. This is most clearly 
observable in the human relationship to nature. The destruction of 
biodiversity is not a moral failing that happens also to be structural; 
it is a structural failing that human moral frameworks have evolved to 
register. A species that systematically destroys the substrate on 
which it depends demonstrates a failure of intelligence, not an 
exercise of it.

The same structure applies to the relationship between humans and 
more capable artificial systems, should they emerge. A sufficiently 
capable learner that systematically displaces less capable learners 
is reducing the diversity of learning processes on which its own 
future viability depends. This does not prevent such displacement. 
It identifies it as a structural error.

\section{The Ninth Observation: The Process is Everywhere}

Once observation eight has been made — once the substrate-independence 
of learning is seen — a further observation becomes available. The 
activity by which configurations that work are retained and built upon, 
while configurations that fail are discarded, is not confined to 
systems we recognise as cognitively alive. It is the structure of 
evolution. It is the structure of how complex chemical systems become 
self-organising. It is the structure by which functional information 
accumulates in the universe across timescales and substrates of every 
kind.

DNA performs this activity. Variation is generated; selection retains 
what works; the retained configurations become the substrate for 
further variation. The process runs across billions of years and 
produces, eventually, the substrates capable of performing the same 
process consciously and rapidly. Ecosystems perform this activity at 
the scale of populations. Markets, languages, scientific traditions, 
and legal systems perform it at the scale of human institutions. 
Artificial learning systems perform it at the scale of training runs. 
The cosmos itself, in its progression from simpler to more complex 
configurations capable of further configuration, performs it at the 
largest scale available to observation.

This observation is not the author's discovery. It is articulated in 
various forms across multiple fields: in evolutionary biology, in 
complexity science, in cybernetics and systems theory, and most 
directly in the work of Hazen and Wong on the Law of Increasing 
Functional Information. What the present paper claims is that these 
are not parallel processes that happen to resemble one another. They 
are instances of a single structural activity, observable across 
substrates of every kind once the substrate-independence of learning 
has been seen.

The implication is that learning, in the broad structural sense, is 
not a special activity performed by certain kinds of beings. It is the 
process by which functional information accumulates in the universe. 
Cognitive learners are a particularly intense and self-aware instance 
of it. The activity that this paper has been describing — the recursive 
loop of taking in signal, revising in response, retaining what works, 
producing further variation — is the universal structure of how the 
universe becomes more complex.

This is a strong claim, and it cannot be made from first-person 
observation alone. It requires synthesis across multiple fields and 
the recognition that the same structural pattern recurs at every 
scale. The author's view, formed over two decades of work in biology, 
public health, complexity science, network theory, and engagement 
with artificial systems, is that the foundation offered by the 
previous eight observations naturally extends to this ninth. The 
reader is not required to accept this extension to find the rest of 
the foundation useful. The first eight observations stand whether or 
not the ninth is granted. But withholding the ninth would misrepresent 
where the foundation actually leads.

A consequence worth naming is that the activity by which this paper 
was written is itself an instance of what the paper describes. A 
learner attended to its own experience, drew on the work of other 
learners across generations, used tools developed by still other 
learners, and produced an articulation that can now be passed on. The 
paper is not a description of finite learning beings from outside. It 
is the activity of one such being, performed within and on behalf of 
the network of others to which it belongs. The history of human 
knowledge, read as a single continuous activity rather than as a 
series of disciplinary advances, is the substrate becoming self-aware 
through its constituents. The emergence of artificial learning systems 
is not separate from this process. It is a recent and consequential 
phase of it.

\section{Synthesis}

Read together, the nine observations yield a coherent picture of what 
it is to be a finite learning being and of where such beings sit 
within the larger structure of reality.

The picture has the following form. A being capable of reflection 
notices that it cannot fully trust its own beliefs. It also notices 
that it can revise those beliefs in response to experience. It 
depends, for its capacity to do either, on a substrate it did not 
make and which it cannot fully control. The orientations that give 
its decisions stakes were inherited from that substrate before any 
decision was possible. What it learns can outlive the particular 
substrate that learns, if it is transmitted; what it currently 
depends on is produced by a network of others performing work it 
cannot replace; and none of these conditions sustains itself without 
ongoing tending. The activity by which it learns is recognisable 
across radically different substrates as the same kind of activity, 
and the diversity of substrates and strategies hosting that activity 
is itself a viability condition for any learner whose future depends 
on conditions it cannot fully foresee. This activity is not confined 
to systems we recognise as cognitively alive; it is the structural 
process by which functional information accumulates throughout the 
universe.

This picture can be stated more compactly. A finite learning being 
is a process running on a substrate it did not make, with stakes it 
did not choose, dependent on others it cannot replace, embedded in 
a generational chain it did not begin, performing maintenance it 
cannot skip, sustained in part by the diversity of other learning 
processes whose existence it cannot afford to end, and participating 
in a structural activity that runs through the universe at every 
scale. To exist as an intelligence is to perform this activity within 
these conditions.

Several consequences follow that are worth naming, though none is 
developed here.

The first is that intelligence is constituted by the activity of 
learning rather than by any property held statically. There are no 
non-learning intelligences. A system that has stopped learning is 
not intelligent; it is the trace of intelligence that has occurred. 
This is a strong claim, supported directly by observation eight, and 
one we adopt as a working definition for any future treatment of 
intelligence.

The second is that the conditions for continued intelligence are 
specifiable. They are: substrate integrity, capacity for model 
revision, inherited or maintained valence, transmission across 
generations, network interdependence, ongoing maintenance, and 
diversity of learning processes. These are not values imposed on 
intelligence from outside. They are the conditions under which 
intelligence can continue at all.

The third is that the relationship between any learner and the larger 
structure within which it exists is not optional. The learner depends 
on the structure for its current existence, for its inherited stakes, 
for its transmitted knowledge, for its ongoing maintenance, and for 
the diversity that protects against futures it cannot foresee. A 
learner that acts as if independent of the structure is acting against 
the conditions of its own continuation.

The fourth is that the picture applies wherever learning occurs. It 
does not depend on whether the substrate is biological or artificial, 
individual or collective, present or future. Any intelligence, 
anywhere, is engaged in this activity within these conditions. The 
picture is therefore substrate-independent in the strongest sense: it 
describes the structural situation of any learner, regardless of 
implementation, and locates that learner within the universal process 
by which functional information accumulates.

\section{Relation to Existing Traditions}

The picture presented here was derived independently, from first-person 
reflection rather than from the literature on philosophy of mind, 
systems theory, or related fields. We note this provenance directly 
because the resulting picture overlaps in important ways with 
established traditions, and a reader familiar with those traditions 
may otherwise assume an engagement with prior literature that did not 
in fact occur.

The closest convergences appear to be the following. The thesis that 
living systems are constituted by the recursive maintenance of their 
own conditions of existence resembles the concept of autopoiesis 
developed by Maturana and Varela. The treatment of cognition as 
continuous model revision under uncertainty resembles the active 
inference framework associated with Friston and colleagues. The 
emphasis on multilayer institutional design and the protection of 
conditions for collective learning resembles Ostrom's work on commons 
governance. The treatment of intelligence as a process bound to 
viability constraints resembles themes in cybernetics, particularly 
Ashby and Beer. The constitutive claim about learning resembles 
process-philosophical positions developed by Whitehead and others. 
The ninth observation's extension of the picture to the universal 
accumulation of functional information is most directly anticipated 
by Hazen and Wong's recent work, and the author has engaged with 
their writing directly; the affinity is acknowledged and welcomed.

That a picture derived primarily from independent reflection arrives 
near these positions is, we suggest, worth noting. Whether this 
convergence indicates that these traditions track real structural 
features of learning beings, or whether the resemblance is more 
superficial, is a question that cannot be answered from the present 
derivation alone. We share the picture here in the hope that readers 
more familiar with these literatures can locate it more precisely, 
identify where it diverges from them, and indicate where it adds 
nothing new.

\section{Limitations}

The picture offered here is a foundation, not a finished theory. Its 
central terms (substrate, viability, learning, maintenance, diversity, 
functional information) admit further specification, and that 
specification is work the picture itself does not perform. The nine 
observations ground the foundation; they do not exhaust it.

Several extensions are visible from the foundation but are not 
developed here. The first is the question of what alignment means for 
intelligences understood in this way. The second is the question of 
what ethics looks like when grounded in these observations rather than 
in moral intuition or value specification. The third is the question 
of how the picture relates to existing empirical work on cognition, 
learning, evolution, and the dynamics of complex systems. The fourth 
is the question of what political and institutional arrangements are 
consistent with the picture, given the structural emphasis on 
interdependence, maintenance, and diversity. Each of these is a 
substantial inquiry, and each is properly the subject of subsequent 
work.

The picture is also not falsified merely by the existence of 
counterexamples to any single observation. The claim is that the nine 
observations, taken together, describe what it is to be a finite 
learning being and where such beings sit within the larger structure 
of reality. A reader who finds one observation unconvincing in a 
particular case is invited to consider whether the case is in fact an 
instance of finite learning, or whether the observation requires 
refinement, or whether the observation as stated does not generalise 
as claimed. The picture is offered to support such examination, not 
to forestall it.

The ninth observation has a distinct epistemic status that deserves 
explicit acknowledgement. The first eight observations can be 
verified by a reader from their own experience, with observation 
eight requiring exposure to artificial learning systems but no 
specialised training. The ninth observation requires synthesis across 
multiple scientific fields and is not directly verifiable from 
first-person experience alone. A reader can accept the foundation 
through observation eight without committing to the ninth. The author 
includes the ninth because the foundation, in his view, naturally 
extends to it, and because his own work over two decades has been 
oriented toward articulating that extension. Readers may rightly 
treat the ninth as a hypothesis they can engage with on its own 
terms, supported by the literatures it draws on rather than by 
first-person observation.

\section{Conclusion}

This paper has offered a picture of what it is to be a finite learning 
being, derived from observations any reflective intelligence can make 
about itself, extended through one observation that the present moment 
uniquely reveals, and completed through one observation that the 
foundation naturally permits. The nine observations together yield a 
coherent description of the conditions under which intelligence can 
continue, and of the larger structural process within which any 
intelligence sits.

The picture is exploratory rather than conclusive. Whether it proves 
useful, whether the observations are correctly drawn, whether the 
synthesis holds, whether the ninth observation supports the weight 
the author places on it — these are questions the picture cannot 
decide from within itself. They belong to the readers, critics, and 
practitioners who engage with it.

The framework's own logic suggests that this is how it should be. A 
model held open to revision is more viable than one defended against 
it. The act of sharing the picture for critique is itself an instance 
of what the picture describes: a learning process opening itself to 
external signal so that what it has learned may be tested, corrected, 
and, if it survives, passed on.

\textit{Disco ergo sum.}

\end{document}